\section{Problem}
\label{sec:problem}

A \emph{run} is one execution of a task on one set of inputs, in one pod. Before the pod starts, an approach
may look at three things. The first is the task's code as it was submitted: \code{job.py} and the Job manifest
for a raw Job, and the task module for a Flyte task (\cref{sec:twins}). The second is the metadata a scheduler
can read about this run's inputs: their sizes and formats and whatever the format's header exposes, such as
the row count and schema of a Parquet file or the shape and dtype of a \code{.npy} array. The third is what
earlier runs of the same task used: their peak memory, peak and mean cores, peak GPU memory and utilization,
wall time and input sizes. These are reported the way the platform's monitoring would show them, which means
the container's own peak for a Job and the peak including the runtime for a Flyte task.

An approach is never given the task's family, the hidden settings it was generated with, the corpus it came
from, or anything measured on the run itself. A unit test checks that changing those fields leaves the
features and the prompt unchanged. The code does identify the family in practice: each task is named after its
family, and each family's generator has a fixed function name and input schema. \Cref{sec:res-manifest}
re-scores with neutral names.

From what it can see, the approach chooses a request on three axes (\cref{tab:grid}). A GPU request lands on
the cheapest machine that can hold it (\cref{tab:ladder}). On a GPU node the host memory and cores come with
the instance, so they are not scored.

\begin{table}[tbp]
\centering\small
\caption{The request grid.}
\label{tab:grid}
\begin{tabular}{@{}lrl@{}}
\toprule
axis & cells & values \\
\midrule
memory & 17 & 256, 384, 512, 768\,Mi; 1, 1.5, 2, 3, 4, 6, 8, 12, 16, 24, 32, 48, 64\,Gi (half powers of two) \\
CPU    & 16 & 0.25, 0.5, 0.75, 1, 1.25, 1.5, 1.75, 2, 2.5, 3, 3.5, 4, 5, 6, 7, 8 cores \\
GPU    & 18 & none, or peak GPU memory 1, 1.5, 2, 3, 4, 6, 8, 10, 12, 15, 22, 32, 44, 64, 88, 120, 176\,GiB \\
\bottomrule
\end{tabular}
\end{table}

\begin{table}[tbp]
\centering\small
\caption{The GPU machine ladder: the cluster's GPU node pools, with on-demand \code{us-east-1} prices for the
whole node. Usable GiB is NVML's total, rounded down. The A10G is never the cheapest shape for any cell.}
\label{tab:ladder}
\begin{tabular}{@{}llrrr@{}}
\toprule
shape & instance & GPUs & usable GiB & \$/h \\
\midrule
T4      & g4dn.xlarge  & 1 & 15         & 0.526 \\
L4      & g6.2xlarge   & 1 & 22         & 0.978 \\
A10G    & g5.2xlarge   & 1 & 22         & 1.212 \\
L40S    & g6e.2xlarge  & 1 & 44         & 2.242 \\
4$\times$L4   & g6.12xlarge  & 4 & 4$\times$22 & 4.602 \\
8$\times$T4   & g4dn.metal   & 8 & 8$\times$15 & 7.824 \\
4$\times$L40S & g6e.12xlarge & 4 & 4$\times$44 & 10.493 \\
\bottomrule
\end{tabular}
\end{table}
